%------------------------------------------------------------------------------%
\section{Equações Trigonométricas}

Uma equação trigonométrica é uma equação que envolve funções trigonométricas
como seno, cosseno, tangente, cotangente, secante ou cossecante de uma variável
trigonométrica. O objetivo da resolução de equações trigonométricas é encontrar
todas as soluções possíveis para a equação.

As equações abaixo são exemplos de equações trigonométricas:
\begin{gather*}
  \sin(x) = \frac{1}{2} \\[2mm]
  \cos^2(x)-\cos(x) = 0 \\[3mm]
  \cos(x)\left(\tan(x)-1\right) = 0
\end{gather*}

Para resolvermos uma equação trigonométrica, precisamos nos atentar em qual
intervalo é pedida a solução.

A resolução de equações trigonométricas pode ser feita por meio de diversas
técnicas, dependendo da complexidade da equação. Alguns dos métodos mais comuns
incluem:
\begin{enumerate}
  \item Equações simples: são aquelas que envolvem apenas uma função
    trigonométrica. Nesse caso, é possível usar propriedades básicas das
    funções trigonométricas para encontrar as soluções.

  \item Equações quadráticas: são aquelas que envolvem o quadrado de uma função
    trigonométrica. Nesse caso, é possível fazer uma substituição adequada para
    transformar a equação em uma equação quadrática comum, que pode ser
    resolvida usando as técnicas usuais.

  \item Equações com produto nulo: são aquelas em que uma ou mais funções
    trigonométricas multiplicadas entre si resultam em zero. Nesse caso, é
    possível usar a propriedade do produto nulo para encontrar as soluções.

  \item Uso de identidades: muitas vezes, as equações trigonométricas podem ser
    resolvidas por meio da aplicação de identidades trigonométricas, como as
    identidades de soma e diferença, identidades de produto e identidades de
    duplo ângulo.
\end{enumerate}

Em geral, a resolução de equações trigonométricas requer um bom conhecimento
das propriedades e identidades das funções trigonométricas, bem como das
técnicas algébricas usadas para resolver equações. Vamos resolver alguns
exemplos.

%------------------------------------------------------------------------------%
\begin{example}
  Resolva a seguinte equação trigonométrica \(\sin(\theta) = \frac{\sqrt{2}}{2}\).

  \solution
  \emph{No intervalo $[0, 2\pi]$}

  Note que, nesse caso, queremos encontrar os arcos no intervalo
  $[0, 2\pi]$, cujo seno vale $\sfrac{\sqrt{2}}{2}$. Sabemos que
  $\sin(\theta)$ é positivo no primeiro e segundo quadrante. No primeiro
  quadrante o ângulo é $\sfrac{\pi}{4}$ (usamos sempre o arco do 1\textsuperscript{o}
  como referência). Para encontrar o arco no segundo quadrante, usamos a
  expressão
  \[
    \pi-\frac{\pi}{4} = \frac{3\pi}{4},
  \]
  conforme mostra a figura.

  \IncludeGraphicCentredEx{equa1png}[0.5][-1][-1]

  Portanto, a solução da equação pedida é
  \[
    S = \left\{ \frac{\pi}{4}, \frac{3\pi}{4} \right\}.
  \]

  \emph{Solução geral}

  A solução geral é dada por:
  \[
    S =
    \left\{
       x \in \R | x=\frac{\pi}{4}+2k\pi
      \; \text{ ou } \;
      x = \frac{3\pi}{4} +2k\pi, \quad k \in \Z
    \right\}.
  \]
\end{example}
%------------------------------------------------------------------------------%

%------------------------------------------------------------------------------%
\begin{example}
  Resolva a equação
  \(2\cos\left(\frac{\theta}{3}\right) - \sqrt{3} = 0\),
  onde $0 \leq \theta \leq 3\pi$.

  \solution
  A equação
  \[
    2\cos\left(\frac{\theta}{3}\right) - \sqrt{3} = 0
  \]
  é equivalente a
  \[
    \cos\left(\frac{\theta}{3}\right) = \frac{\sqrt{3}}{2}.
  \]
  Como $0 \leq \theta \leq 3\pi$, então,
  \[
    0 \leq \frac{\theta}{3} \leq \pi.
  \]
  Uma vez que o único arco no intervalo $[0, \pi]$,
  que possui cosseno igual a
  \[
    \sfrac{\sqrt{3}}{2}$ é $\sfrac{\pi}{6},
  \]
  temos que
  \begin{align*}
    \frac{\theta}{3} & = \frac{\pi}{6}  \\
    \theta           & = \frac{\pi}{2}.
  \end{align*}
  Portanto,
  \[
    S = \left\{\frac{\pi}{2} \right\}.
  \]
\end{example}
%------------------------------------------------------------------------------%

%------------------------------------------------------------------------------%
\begin{example}
  Resolva a equação $2\cos(x)= -1$ no intervalo $[0,2\pi]$.

  \solution
  Primeiro, vamos isolar o $\cos(x)$ dividindo ambos os lados por 2:
  \[
    \cos(x)= -\frac{1}{2}
  \]

  Sabemos que o cosseno é negativo no segundo e terceiro quadrante, para
  encontrarmos quais são esses arcos, pegamos como referência o arco no
  primeiro quadrante que tem cosseno igual a
  $\sfrac{1}{2}$, isto é,
  $\sfrac{\pi}{3}$.

  \IncludeGraphicCentredEx{equa2}[0.5][-1][-1]

  Conforme mostra a figura, o arco no segundo quadrante é
  \[
    \pi-\frac{\pi}{3}= \frac{2\pi}{3}
  \]
  e o arco no terceiro quadrante é
  \[
    \pi+\frac{\pi}{3}= \frac{4\pi}{3}.
  \]
  Assim,
  \[
    \cos(x) = -\frac{1}{2}
  \]
  ocorre quando $x$ é um ângulo de
  \[
    \frac{2\pi}{3}
    \quad\text{ ou }\quad
    \frac{4\pi}{3}
  \]
  (ou seus equivalentes negativos).
  Portanto, a solução da equação no intervalo $[0, 2\pi]$ é
  \[
    S= \left\{\frac{2\pi}{3}, \frac{4\pi}{3} \right\}
  \]
\end{example}
%------------------------------------------------------------------------------%

%------------------------------------------------------------------------------%
\begin{example}
  Resolva a equação
  \[
    \left(\sin(\theta) -\frac{1}{2} \right)
    \left(\tan(\theta) -1 \right) = 0,
  \]
  para $0 \leq \theta < \sfrac{\pi}{2}$.

  \solution
  Essa equação é do tipo produto nulo. Sabemos que produto $ab$ só é zero se
  $a = 0$ ou $ b = 0$.
  Portanto:
  \[
    \sin(\theta) -\frac{1}{2} = 0
    \quad\text{ ou }\quad
    \tan(\theta) -1 = 0
  \]
  Analisando cada caso, separadamente
  \begin{align*}
    \sin(\theta) -\frac{1}{2} & =0              \\
    \sin(\theta)              & =\frac{1}{2}    \\
    \theta                    & = \frac{\pi}{6}
  \end{align*}
  ou
  \begin{align*}
    \tan(\theta) -1 & = 0            \\
    \tan(\theta)    & = 1            \\
    \theta          & =\frac{\pi}{4}
  \end{align*}
  Assim, a solução da equação pedida é:
  \[
    S = \left\{ \frac{\pi}{6}, \frac{\pi}{4} \right\}.
  \]
\end{example}
%------------------------------------------------------------------------------%

%------------------------------------------------------------------------------%
\begin{example}
  Resolva a equação  $2\sin^2 x - \sin x - 1 = 0$ no intervalo $[0, 2\pi]$.

  \solution
  Neste caso, temos uma equação quadrática trigonométrica que pode ser
  resolvida através da mudança de variável $y = \sin(x)$. Dessa forma, obtemos
  a equação do segundo grau
  \[
    2y^2  - y - 1 = 0,
  \]
  cujas raízes são
  \[
    y = 1
    \quad\text{ e }\quad
    y = -\frac{1}{2}.
  \]
  Portanto, podemos encontrar os valores de $x$
  correspondentes da seguinte maneira.

  Para $y = 1$, temos $\sin(x) = 1$, o que implica em
  \(
    x = \frac{\pi}{2}.
  \)

  Para $y = -\sfrac{1}{2}$, temos $\sin(x) = -\sfrac{1}{2}$.
  Isso ocorre em dois
  quadrantes, 3\textsuperscript{o} e 4\textsuperscript{o} quadrantes.
  Portanto, os valores
  possíveis de $x$ são
  \[
    \pi + \frac{\pi}{6} = \frac{7\pi}{6}
    \qquad \text{(3\textsuperscript{o} Q)}
  \]
  ou
  \[
    2\pi - \frac{\pi}{6} = \frac{11\pi}{6}
    \qquad \text{(4\textsuperscript{o} Q)}.
  \]
  Assim o conjunto solução da equação original é
  \[
    S = \left\{ \frac{\pi}{2}, \frac{7\pi}{6}, \frac{11\pi}{6} \right\}
  \]
\end{example}
%------------------------------------------------------------------------------%

